Четвертый раздел посвящен планированию реализации выбранной траектории развития посредством построения графика Ганта. В данном разделе необходимо систематизировать мероприятия индивидуального плана развития, определив их последовательность и продолжительность во времени. Построение графика требует учета текущей академической нагрузки, профессиональной деятельности, а также иных значимых жизненных ролей для обеспечения сбалансированного распределения временных ресурсов.

Из таблиц индивидуального плана развития (\cref{tab:idp,tab:idp_freelance}) отобраны ключевые мероприятия для реализации в течение ближайшего полугодия. График Ганта представлен в \cref{tab:gantt}.

\begin{figure}[H]
    \centering
    \includegraphics[width=\textwidth]{figs/gant.pdf}
    \caption{График Ганта}
    \label{tab:gantt}
% -- XE Latex Compiler
% \documentclass[tikz]{standalone}
% \usepackage{pgfgantt}

% % Настройки для русского языка
% \usepackage{fontspec}
% \setmainfont{Arial} % Или любой другой системный шрифт с кириллицей

% \begin{document}
% \begin{ganttchart}[
%     hgrid,
%     vgrid,
%     x unit=2.5cm,
%     y unit title=0.7cm,
%     y unit chart=0.7cm,
%     title label font=\itshape\footnotesize,
%     bar label font=\footnotesize,
%     bar/.append style={fill=gray!40, draw=black},
%     bar height=0.5
% ]{1}{6}
%     \gantttitle{График индивидуального плана развития}{6} \\
%     \gantttitle{Месяц 1}{1}
%     \gantttitle{Месяц 2}{1}
%     \gantttitle{Месяц 3}{1}
%     \gantttitle{Месяц 4}{1}
%     \gantttitle{Месяц 5}{1}
%     \gantttitle{Месяц 6}{1} \\
    
%     \ganttbar{Обучение в университете}{1}{6} \\
%     \ganttbar{Изучение языков МЭК 61131-3}{1}{2} \\
%     \ganttbar{Освоение среды TIA Portal v16}{2}{4} \\
%     \ganttbar{Развитие навыков управления проектами}{3}{4} \\
%     \ganttbar{Изучение SCADA-систем и HMI}{4}{5} \\
%     \ganttbar{Изучение ГОСТ по оформлению АСУ ТП}{5}{5} \\
%     \ganttbar{Технический английский язык}{1}{6} \\
% \end{ganttchart}
% \end{document}
\end{figure}